\section{結果}

各条件の平均実行時間を表~\ref{tab:timing}に示す．入力サイズ8,000件では，バブルソートの平均は\SI{194.74}{ms}，クイックソートの平均は\SI{2.372}{ms}であり，両者の比は82.1であった．

\begin{table}[tbp]
  \centering
  \caption{入力サイズ別の平均実行時間（各5回の架空データ）}
  \label{tab:timing}
  \begin{tabular}{rrrr}
\toprule
入力サイズ $n$ & バブルソート（ms） & クイックソート（ms） & 比率 \\
\midrule
1,000 & 3.12 & 0.234 & 13.3 \\
2,000 & 12.24 & 0.513 & 23.9 \\
4,000 & 48.76 & 1.116 & 43.7 \\
8,000 & 194.74 & 2.372 & 82.1 \\
\bottomrule
\end{tabular}

\end{table}

入力サイズと平均実行時間の関係を図~\ref{fig:timing}に示す．縦軸は両系列を同時に比較できるよう対数目盛とした．全条件でバブルソートの平均実行時間がクイックソートを上回り，入力サイズの増加とともに差が広がった．

\begin{figure}[tbp]
  \centering
  \begin{tikzpicture}
    \begin{axis}[
      width=0.82\textwidth,
      height=62mm,
      xlabel={入力サイズ $n$（千件）},
      ylabel={平均実行時間（ms）},
      ymode=log,
      log basis y=10,
      xtick={1,2,4,8},
      grid=major,
      grid style={black!15},
      legend style={draw=none,fill=none,at={(0.03,0.97)},anchor=north west},
      tick label style={font=\small},
      label style={font=\small},
    ]
      \addplot[black,thick,mark=*] table[x=n_thousands,y=mean_ms]{generated/bubble.dat};
      \addlegendentry{バブルソート}
      \addplot[black!60,thick,dashed,mark=square*] table[x=n_thousands,y=mean_ms]{generated/quick.dat};
      \addlegendentry{クイックソート}
    \end{axis}
  \end{tikzpicture}
  \caption{入力サイズと平均実行時間の関係}
  \label{fig:timing}
\end{figure}
